\chapter{Interactive Topology Viewer}
\label{ch:viewer}

\section{Purpose and Architecture}
\label{sec:viewer-purpose}

Static PNG renderings of a circuit's DAG (\code{circuit\_graph\_\{pre,post\}\_optimization.png} artifacts of
\cref{tab:output-contract}) are adequate for small circuits but become
difficult to read once a circuit reaches the size of the larger benchmark
circuits (upwards of a hundred gates) or once a reader wants to compare the
pre- and post-optimization state of the same circuit interactively rather
than by flipping between two static images. The interactive viewer
addresses this: every analysis run exports \code{circuit\_topology.json}, a
versioned, self-contained graph description --- one shared topology
(primary inputs, gates, primary outputs, and named nets) with two named
overlays, \code{original} and \code{canonical\_optimized}, each supplying
per-gate cell assignment, load capacitance, rise/fall delay, and ranked
critical paths for that state of the circuit. A small local server,
launched with

\begin{lstlisting}[language=bash, style=pythonstyle, numbers=none]
vlsi-sta view outputs/c15_high_fanout_sizing_fabric
\end{lstlisting}

serves this JSON contract to a browser-based front end and opens it
automatically (suppressible with \code{--no-browser} for remote or scripted
use, with \code{--host}/\code{--port} available for remote access).

\section{Features}
\label{sec:viewer-features}

\begin{itemize}
    \item Mouse- and touch-based panning, cursor-anchored wheel zoom, and a
          fit-to-view command, with keyboard shortcuts (\code{+}, \code{-},
          \code{F}, \code{0}, \code{Escape}) mirroring the on-screen
          controls.
    \item Optional net labels and node selection, with connected-net
          isolation to declutter the view of a densely fanned-out node.
    \item Independent visibility toggles for the \code{original} and
          \code{canonical\_optimized} overlays, so the two states can be
          viewed separately or superimposed.
    \item Critical-path coloring consistent with the static PNG renderings.
    \item \textbf{Dual-state gate inspection}: selecting a gate shows its
          cell, load, and delay in both the original and optimized state
          side by side.
    \item \textbf{Target-size highlighting} for any gate the optimizer
          resized: an amber halo marks a gate resized to an \code{X2}
          target, a wider purple halo marks a gate resized to an \code{X4}
          target. This highlighting is shown only while both overlays are
          simultaneously enabled, since it is inherently a comparison
          between the two states.
      \item \textbf{Analysis and Schematic View modes}: Analysis view is the
          normal mode (similar to the PNG files), but a \emph{Schematic View} feature
          enables the user to see the real circuit elements instead of boxes.
\end{itemize}

\section{Screenshots}
\label{sec:viewer-screenshots}
\vspace{-10pt}
\begin{figure}[H]
    \centering
    \includegraphics[width=\textwidth]{viewer/analysis view}
    \caption{Interactive topology viewer: Analysis View mode.}
\end{figure}

\begin{figure}[H]
    \centering
    \includegraphics[width=\textwidth]{viewer/schematic view}
    \caption{Interactive topology viewer: Schematic View mode.}
\end{figure}
\vspace{-60pt}

\begin{figure}[H]
    \centering
    \includegraphics[width=\textwidth]{viewer/gate select}
    \caption{Interactive topology viewer: gate inspection mode (clicked on a gate).}
\end{figure}

\begin{figure}[H]
    \centering
    \includegraphics[width=\textwidth]{viewer/big circuit}
    \caption{Interactive topology viewer: Schematic view for a larger circuit.}
\end{figure}
